\section{Introduction}\label{sec:intro}

Mechanistic interpretability makes two kinds of causal claim about in-context computation: where a model stores and reads a piece of state, and which internal component carries the effect of an intervention \citep[e.g.,][]{vig2020causal,wang2022ioi,geiger2024finding}. Both are usually measured under a single readout. One such readout is a multiple-choice prompt, which is common in LM evaluation benchmarks and is used in the MCQA task of the MIB interpretability benchmark \citep{mueller2025mib}; circuit analyses have also studied multiple-choice answering directly \citep{lieberum2023chinchilla,wiegreffe2025answer}. Model behaviour is known to be sensitive to prompt format \citep{sclar2023quantifying}, and recent work examines whether steering and function vectors act invariantly across formats and answer encodings \citep{opielka2026causality,gao2026encodings}. Whether the attribution of an intervention's effect to a specific internal component, such as one attention channel at one token, survives a change of readout has received less direct attention.

A concrete case comes from \citet{anonymous2026fitted}. They fit rank-16 activation patches on belief-tracking stories to implement a remapping of locations, and compare them with unfitted PCA patches of the same rank. Exchanging only the attention keys at the critical token, the token that names the location in the critical (move or tell) event, between the learned and the PCA run substantially transfers the learned patch's downstream behaviour. They conclude that the critical-token keys induced by the learned patch contribute causally to the behavioural difference between the learned remap and the PCA patch, and present this as a component-level explanation of the tested intervention effects. Their evaluation lists the candidate locations after the question.

A later token can see a context token only through two quantities: the token's key, which decides how much the later token attends to it, and its value, which decides what is copied when it does. A key therefore acts only through the attention that later queries pay to it. This suggests a prediction: the key channel should matter when later tokens have something to match, such as re-mentioned candidate answers, and should matter little when they do not. For free-form answers, a lookback account of belief tracking places an address in the key of a recalled token, matched by a later query, with attention then moving a payload from that token \citep{prakash2025lookbacks}; a rebinding account of entity tracking finds binding information in query/key subspaces or mainly in key vectors \citep{oh2026rebinding}.

We test this prediction directly. On the synthetic stories of \citet{anonymous2026fitted}, we clamp the key and the value of the single state token separately to those of a source story, at every layer. Clamping both reproduces the source run exactly, so the two channels account for the full effect. We vary only the readout, across five formats: the options-after-question format of \citet{anonymous2026fitted} (\fmtOpt{}), lettered options (\fmtLetter{}), a neutral sentence re-mentioning all six locations (\fmtPost{}), no candidates (\fmtNone{}), and options before the story (\fmtBefore{}), whose listed words cannot attend to the state token under the causal mask. We run \nModels{} instruction-tuned models from 1.5B to 72B parameters in four families, committing predictions before each of four confirmatory runs and scoring scripts before inspecting outputs; the initial 1.5B factorial, the 1.5B localisation and the neutral re-mention effect from 7B upward are exploratory.

The readout decides which channel carries the state (\cref{fig:formats}). When the candidates are listed after the state token (\fmtOpt{}, \fmtLetter{}), the keys carry the location's identity in all \nModels{} models: the key share of the \fmtOpt{} effect ranges from \shareOlmoSeven{} (OLMo-2-7B) to \shareMistralSeven{} (Mistral-7B) in the eight models of 7B and above, and is \shareQwenOnefive{} at Qwen2.5-1.5B. Exact row-restricted swaps at Qwen2.5-1.5B show that the option words themselves read the key: they recover \locOptChoicewords{} (\fmtOpt{}) and \locLetterChoicewords{} (\fmtLetter{}) of the key effect, while question rows recover at most \locLetterQuestion{}. Without re-mention the state is copied from the value, and keys carry at most \noneMax{} nats of identity; with options before the story they carry at most \beforeMax{}. From 7B upward a neutral re-mention also opens the key channel, at an intermediate strength; we did not predict this, and at 1.5B and 3B it carries no positive key identity. The consequence for \citet{anonymous2026fitted} is direct (\cref{fig:frames}). Their learned remap shifts behaviour by nearly the same fraction in every format ($\varphi$ between \phiNoneQwen{} and \phiBeforeQwen{} at Qwen2.5-72B, \phiOptMistral{}--\phiLetterMistral{} at Mistral-24B), but adding its keys to the PCA run recovers \psiOptQwen{} of the learned--PCA difference under \fmtOpt{}, \psiLetterQwen{} under \fmtLetter{}, \psiNoneQwen{} under \fmtNone{}, and $\psiBeforeQwen$ under \fmtBefore{}. The edit is written into both channels, and which one appears to carry it depends on which one the readout reads. \todo{Stage 3: generality to the paint and schedule tasks and localisation at 7--14B.}

Not every prediction held. A preregistered transfer law predicted the free-form effect of the remap from its value-channel completeness $\kappa_V$, which we estimated beforehand from the fixed-value key-exchange data of \citet{anonymous2026fitted}. It fails at Mistral-24B (predicted 0.53 $\pm$ 0.15, observed \phiNoneMistral{}) and is met only at the edge of its margin at Qwen2.5-72B (predicted 0.73 $\pm$ 0.15, observed \phiNoneQwen{}). Its companion prediction, that the remap transfers less to free-form answers than to the multiple-choice format it was fit in, also fails at Mistral-24B (difference $-0.003$ [$-0.019$, $+0.013$]) and holds only weakly at Qwen2.5-72B ($+0.034$ [$+0.018$, $+0.049$]). A preregistered key-share threshold at 24--72B is met only at Qwen2.5-32B; within Qwen2.5 the key share rises from \shareQwenOnefive{} at 1.5B to \shareQwenFourteen{} at 14B and then falls to \shareQwenSeventytwo{} at 72B. A narrower ``listed options only'' account, adopted post hoc at 1.5B, passed a preregistered fresh-seed test in the same model but does not generalise to models of 7B and above, where neutral re-mentions also open the key channel and, with no re-mention, keys still carry up to \noneMax{} nats of identity, outside the $\pm 1$ nat equivalence margin set at 1.5B.

\paragraph{Contributions.}
\begin{itemize}
  \item \textbf{Mechanism.} Using exact key/value clamps at a single state token, we show that answer options listed after the state token, and from 7B upward also a neutral re-mention of the candidates, look in-context state up through the state token's key, while without re-mention the state is copied from its value. Row-restricted splices localise the read to the listed option words at 1.5B (\cref{sec:results}).
  \item \textbf{Scale and controls.} The pattern holds in \nModels{} models from 1.5B to 72B in four families, with a structural control (options before the story) under which keys carry at most \beforeMax{} nats of location identity in every model (\cref{tab:natural}, \cref{fig:scale}).
  \item \textbf{Readout-relativity of an intervention attribution.} After reproducing \citet{anonymous2026fitted} item by item (argmax agreement 0.994--1.000), we show that their learned remap shifts behaviour by nearly the same fraction under every readout, while its attribution to critical-token keys is strong only when the candidates are listed after the state token ($\psi$ \psiOptQwen{}--\psiLetterQwen{} at Qwen2.5-72B), weak under a neutral re-mention (\psiPostQwen{}) and near zero otherwise (\cref{tab:frames}, \cref{fig:frames}).
  \item \textbf{Methodology.} Four preregistrations, with predictions committed before each preregistered run and scoring scripts before inspection; every outcome scored so far is reported, including failed predictions (\cref{app:prereg}) \todo{update when preregistration D (stage 3) is scored}; and code for exact clamps and row-restricted splices, to be released with the paper.
\end{itemize}
